% Part: first-order-logic
% Chapter: sequent-calculus
% Section: provability-consistency

\documentclass[../../../include/open-logic-section]{subfiles}

\begin{document}

\iftag{FOL}
      {\olfileid[hi]{fol}{seq}{prv}}
      {\olfileid[hi]{pl}{seq}{prv}}

\olsection{\usetoken{S}{derivability} और संगतता}

अब हम !!{derivability} संबंध के अनेक गुण स्थापित करेंगे। ये अपने-आप में
रोचक हैं, पर इनमें से प्रत्येक पूर्णता प्रमेय के प्रमाण में भूमिका निभाएगा।

\begin{prop}\ollabel{prop:provability-contr}
  यदि $\Gamma \Proves !A$ और $\Gamma \cup \{!A\}$ असंगत है, तो
  $\Gamma$ असंगत है।
\end{prop}

\begin{proof}
ऐसे परिमित समुच्चय $\Gamma_0$ और $\Gamma_1 \subseteq \Gamma$ हैं कि
$\Log{LK}$, $\Gamma_0 \Sequent !A$ और $!A, \Gamma_1
\Sequent \quad$ को !!{derive} करता है। $\Log{LK}$ में $\Gamma_0 \Sequent
!A$ की !!{derivation} को~$\pi_0$ और $\Log{LK}$ में $\Gamma_1, !A
\Sequent \quad$ की !!{derivation} को~$\pi_1$ मानें। तब हम निम्न अनुक्रम को
!!{derive} कर सकते हैं:
\begin{prooftree}
\AxiomC{}
\RightLabel{$\pi_0$}
\Deduce$ \Gamma_0 \fCenter !A $
\AxiomC{}
\RightLabel{$\pi_1$}
\Deduce$!A, \Gamma_1 \fCenter $
\RightLabel{\Cut}
\BinaryInf$ \Gamma_0,\Gamma_1 \fCenter $
\end{prooftree}

चूँकि $\Gamma_0 \subseteq \Gamma$ और $\Gamma_1 \subseteq \Gamma$, इसलिए
$\Gamma_0 \cup \Gamma_1 \subseteq \Gamma$; अतः $\Gamma$ असंगत है।
\end{proof}

\begin{prop}
\ollabel{prop:prov-incons}
$\Gamma \Proves !A$ तभी और केवल तभी जब $\Gamma \cup \{\lnot !A\}$ असंगत हो।
\end{prop}

\begin{proof}
पहले मान लें कि $\Gamma \Proves !A$; अर्थात् कोई
!!a{derivation}~$\pi_0$ ऐसी है जिसका अंतिम अनुक्रम $\Gamma \Sequent !A$ है।
$\LeftR{\lnot}$ नियम जोड़ने पर $\lnot !A, \Gamma
\Sequent \quad$ की !!a{derivation} मिलती है; अर्थात्
$\Gamma \cup \{\lnot !A\}$ असंगत है।

यदि $\Gamma \cup \{\lnot !A\}$ असंगत है, तो कोई
!!a{derivation}~$\pi_1$ ऐसी है जिसका अंतिम अनुक्रम $\lnot !A, \Gamma \Sequent \quad$
है। निम्नलिखित $\Gamma \Sequent !A$ की !!a{derivation} है:
\begin{prooftree}
  \Axiom$!A \fCenter !A$
  \RightLabel{\RightR{\lnot}}
  \UnaryInf$\fCenter !A, \lnot !A$
  \AxiomC{}
  \RightLabel{$\pi_1$}
  \Deduce$\lnot !A, \Gamma \fCenter$
  \RightLabel{\Cut}
  \BinaryInf$\Gamma \fCenter !A$
\end{prooftree}
\end{proof}

\begin{prob}
सिद्ध कीजिए कि $\Gamma \Proves \lnot !A$ तभी और केवल तभी जब
$\Gamma \cup \{!A\}$ असंगत हो।
\end{prob}

\begin{prop}\ollabel{prop:explicit-inc}
  यदि $\Gamma \Proves !A$ और $\lnot !A \in \Gamma$, तो $\Gamma$ असंगत है।
\end{prop}

\begin{proof}
  मान लें कि $\Gamma \Proves !A$ और $\lnot !A \in \Gamma$। तब कोई
  !!a{derivation}~$\pi$ ऐसी है जिसका अंतिम अनुक्रम $\Gamma_0 \Sequent !A$ है। अनुक्रम
  $\lnot !A, \Gamma_0 \Sequent \quad$ भी !!{derivable} है:
  \begin{prooftree}
    \AxiomC{}
    \RightLabel{$\pi$}
    \Deduce$\Gamma_0 \fCenter !A$
    \Axiom$!A \fCenter !A$
    \RightLabel{\LeftR{\lnot}}
    \UnaryInf$\lnot !A, !A \fCenter$
    \RightLabel{\LeftR{\Exchange}}
    \UnaryInf$!A, \lnot !A \fCenter$
    \RightLabel{\Cut}
    \BinaryInf$\Gamma_0, \lnot !A \fCenter$
  \end{prooftree}
  चूँकि $\lnot !A \in \Gamma$ और $\Gamma_0 \subseteq \Gamma$, इससे सिद्ध
  होता है कि $\Gamma$ असंगत है।
\end{proof}

\begin{prop}\ollabel{prop:provability-exhaustive}
  यदि $\Gamma \cup \{!A\}$ और $\Gamma \cup \{\lnot !A\}$ दोनों असंगत हैं,
  तो $\Gamma$ असंगत है।
\end{prop}

\begin{proof}
ऐसे परिमित समुच्चय $\Gamma_0 \subseteq \Gamma$ और $\Gamma_1
\subseteq \Gamma$ तथा $\Log{LK}$-!!{derivation}s $\pi_0$ और $\pi_1$ हैं,
जिनके अंतिम अनुक्रम क्रमशः $!A, \Gamma_0 \Sequent \quad$ और
$\lnot !A, \Gamma_1 \Sequent
\quad$ हैं। तब हम निम्न अनुक्रम को !!{derive} कर सकते हैं:
\begin{prooftree}
\AxiomC{}
\RightLabel{$\pi_0$}
\Deduce$ !A, \Gamma_0 \fCenter $
\RightLabel{\RightR{\lnot}}
\UnaryInf$ \Gamma_0 \fCenter \lnot !A$
\AxiomC{}
\RightLabel{$\pi_1$}
\Deduce$\lnot !A, \Gamma_1 \fCenter  $
\RightLabel{\Cut}
\BinaryInf$ \Gamma_0, \Gamma_1 \fCenter $
\end{prooftree}
चूँकि $\Gamma_0 \subseteq \Gamma$ और $\Gamma_1 \subseteq \Gamma$, इसलिए
$\Gamma_0 \cup \Gamma_1 \subseteq \Gamma$। अतः $\Gamma$ असंगत है।
\end{proof}

\end{document}
